%!TeX root=../pridetop.tex
\chapter[Chapter \thechapter]{}
	
\begin{figure}[t!]
\centering
\includegraphics[width=\linewidth]{37top}
\captionlistentry{His parting obeisance}
\end{figure}


\lettrine[lines=6,image=true]{initials/chap37t}{he} two gentlemen left Rosings the next morning; and Mr Collins having been in waiting near the lodges, to make them his parting obeisance, was able to bring home the pleasing intelligence of their appearing in very good health, and in as tolerable spirits as could be expected, after the melancholy scene so lately gone through at Rosings. To Rosings he then hastened to console Lady Catherine and her daughter; and on his return brought back, with great satisfaction, a message from her Ladyship, importing that she felt herself so dull as to make her very desirous of having them all to dine with her.

Elizabeth could not see Lady Catherine without recollecting that, had she chosen it, she might by this time have been presented to her as her future niece; nor could she think, without a smile, of what her Ladyship's indignation would have been. <What would she have said? how would she have behaved?> were the questions with which she amused herself.

Their first subject was the diminution of the Rosings' party. <I assure you, I feel it exceedingly,> said Lady Catherine; <I believe nobody feels the loss of friends so much as I do. But I am particularly attached to these young men; and know them to be so much attached to me! They were excessively sorry to go! But so they always are. The dear Colonel rallied his spirits tolerably till just at last; but Darcy seemed to feel it most acutely—more, I think, than last year. His attachment to Rosings certainly increases.>

Mr Collins had a compliment and an allusion to throw in here, which were kindly smiled on by the mother and daughter.

Lady Catherine observed, after dinner, that Miss Bennet seemed out of spirits; and immediately accounting for it herself, by supposing that she did not like to go home again so soon, she added,—

<But if that is the case, you must write to your mother to beg that you may stay a little longer. Mrs Collins will be very glad of your company, I am sure.>

<I am much obliged to your Ladyship for your kind invitation,> replied Elizabeth; <but it is not in my power to accept it. I must be in town next Saturday.>

<Why, at that rate, you will have been here only six weeks. I expected you to stay two months. I told Mrs Collins so before you came. There can be no occasion for your going so soon. Mrs Bennet could certainly spare you for another fortnight.>

<But my father cannot. He wrote last week to hurry my return.>

<Oh, your father, of course, may spare you, if your mother can. Daughters are never of so much consequence to a father. And if you will stay another \textit{month} complete, it will be in my power to take one of you as far as London, for I am going there early in June, for a week; and as Dawson does not object to the barouche-box, there will be very good room for one of you—and, indeed, if the weather should happen to be cool, I should not object to taking you both, as you are neither of you large.>

\begin{figure}[tbh]
\centering
\includegraphics[width=.7\linewidth]{37dawson}
\captionlistentry{Dawson}
\end{figure}

<You are all kindness, madam; but I believe we must abide by our original plan.>

Lady Catherine seemed resigned. <Mrs Collins, you must send a servant with them. You know I always speak my mind, and I cannot bear the idea of two young women travelling post by themselves. It is highly improper. You must contrive to send somebody. I have the greatest dislike in the world to that sort of thing. Young women should always be properly guarded and attended, according to their situation in life. When my niece Georgiana went to Ramsgate last summer, I made a point of her having two men-servants go with her. Miss Darcy, the daughter of Mr Darcy of Pemberley, and Lady Anne, could not have appeared with propriety in a different manner. I am excessively attentive to all those things. You must send John with the young ladies, Mrs Collins. I am glad it occurred to me to mention it; for it would really be discreditable to \textit{you} to let them go alone.>

<My uncle is to send a servant for us.>

<Oh! Your uncle! He keeps a man-servant, does he? I am very glad you have somebody who thinks of those things. Where shall you change horses? Oh, Bromley, of course. If you mention my name at the Bell, you will be attended to.>

Lady Catherine had many other questions to ask respecting their journey; and as she did not answer them all herself attention was necessary—which Elizabeth believed to be lucky for her; or, with a mind so occupied, she might have forgotten where she was. Reflection must be reserved for solitary hours: whenever she was alone, she gave way to it as the greatest relief; and not a day went by without a solitary walk, in which she might indulge in all the delight of unpleasant recollections.

Mr Darcy's letter she was in a fair way of soon knowing by heart. She studied every sentence; and her feelings towards its writer were at times widely different. When she remembered the style of his address, she was still full of indignation: but when she considered how unjustly she had condemned and upbraided him, her anger was turned against herself; and his disappointed feelings became the object of compassion. His attachment excited gratitude, his general character respect: but she could not approve him; nor could she for a moment repent her refusal, or feel the slightest inclination ever to see him again. In her own past behaviour, there was a constant source of vexation and regret: and in the unhappy defects of her family, a subject of yet heavier chagrin. They were hopeless of remedy. Her father, contented with laughing at them, would never exert himself to restrain the wild giddiness of his youngest daughters; and her mother, with manners so far from right herself, was entirely insensible of the evil. Elizabeth had frequently united with Jane in an endeavour to check the imprudence of Catherine and Lydia; but while they were supported by their mother's indulgence, what chance could there be of improvement? Catherine, weak-spirited, irritable, and completely under Lydia's guidance, had been always affronted by their advice; and Lydia, self-willed and careless, would scarcely give them a hearing. They were ignorant, idle, and vain. While there was an officer in Meryton, they would flirt with him; and while Meryton was within a walk of Longbourn, they would be going there for ever.

Anxiety on Jane's behalf was another prevailing concern; and Mr Darcy's explanation, by restoring Bingley to all her former good opinion, heightened the sense of what Jane had lost. His affection was proved to have been sincere, and his conduct cleared of all blame, unless any could attach to the implicitness of his confidence in his friend. How grievous then was the thought that, of a situation so desirable in every respect, so replete with advantage, so promising for happiness, Jane had been deprived, by the folly and indecorum of her own family!

When to these recollections was added the development of Wickham's character, it may be easily believed that the happy spirits which had seldom been depressed before were now so much affected as to make it almost impossible for her to appear tolerably cheerful.

Their engagements at Rosings were as frequent during the last week of her stay as they had been at first. The very last evening was spent there; and her Ladyship again inquired minutely into the particulars of their journey, gave them directions as to the best method of packing, and was so urgent on the necessity of placing gowns in the only right way, that Maria thought herself obliged, on her return, to undo all the work of the morning, and pack her trunk afresh.

When they parted, Lady Catherine, with great condescension, wished them a good journey, and invited them to come to Hunsford again next year; and Miss de Bourgh exerted herself so far as to courtesy and hold out her hand to both.